\Conclusion

В результате проведенного исследования достигнута поставленная цель — разработан прототип многоуровневого фреймворка для микроконтроллера MIK32, который обеспечивает абстракцию аппаратных ресурсов и упрощает процесс создания встроенных приложений.

Основные результаты работы:

Проведен сравнительный анализ существующих фреймворков для микроконтроллеров (FreeRTOS, Zephyr, Mbed OS) и выявлены оптимальные архитектурные принципы для отечественного микроконтроллера MIK32 на базе RISC-V. 

Разработана архитектура фреймворка с разделением архитектурно-зависимых (arch), платформо-зависимых (platform) и общих компонентов, что обеспечивает переносимость и масштабируемость решений.

Создана подсистема GPIO, обеспечивающая единообразную работу с встроенными портами и внешними расширителями типа PCA9554. Разработана подсистема I2C, адаптированная для встроенных систем с поддержкой стандартных режимов работы.

Фреймворк успешно апробирован в двух промышленных проектах: системе управления нагревом Heating Unit и системе управления вентиляторами Midplane, что подтвердило его универсальность и практическую применимость.

Практическая значимость:

Разработанный фреймворк сокращает время разработки встроенных систем на базе MIK32 и способствует ускорению внедрения отечественного микроконтроллера в промышленные применения. Созданный инструментарий имеет стратегическое значение для обеспечения технологической независимости в критически важных отраслях экономики.